\chapter{Conclusiones}
\label{chapter:conclusiones}

En este trabajo llevamos a cabo una refactorización integral de la herramienta \gls{feets} para mejorar su desempeño en contextos de grandes volúmenes de datos, como los planteados por relevamientos astronómicos modernos tales como el \gls{vvv} o \gls{macho}. El nuevo modelo de extracción paralelizado con \textit{Dask} eliminó el cuello de botella presente en la implementación secuencial de la versión original, demostrando tiempos de ejecución significativamente más cortos y escalando adecuadamente ante flujos de datos masivos, propios de la astronomía contemporánea.

Durante el proceso de refactorización logramos, además, mejorar la arquitectura del sistema y estandarizar el código con buenas prácticas de ingeniería del \emph{software}, resultando en un mayor \acrlongsp{mi} global y elevando la calidad interna del proyecto. También modernizamos el \emph{stack} tecnológico de la herramienta, siguiendo la adopción de versiones modernas de Python y la integración de la biblioteca \textit{LightCurve}, que nos permitió incorporar nuevos extractores de características eficientes sin comprometer el desempeño de la aplicación.

Aprovechando las mejoras tanto en el rendimiento como en la calidad del \emph{software}, realizamos cambios a la \gls{api} para también extender su funcionalidad. La nueva implementación de \gls{feets} permite procesar múltiples \glspl{lc} de manera simultánea, simplifica el desarrollo de nuevos extractores mediante un \emph{framework} más eficiente, e introduce una estructura de datos optimizada para la visualización y el análisis posterior de las características extraídas.

Como resultado, la nueva versión de \textit{feets} convierte a la herramienta en una solución robusta, rápida y moderna, que no solo supera las limitaciones de su versión anterior, sino que también está preparada para seguir evolucionando ante los desafíos que presenta el análisis de series temporales en la era del \emph{Big Data}.

\section{Trabajo a futuro}

Cuando evaluamos los tiempos de cómputo para distintas estrategias de planificación, consideramos que otra dirección prometedora sería evaluar el comportamiento del sistema utilizando el planificador \texttt{distributed} que ofrece \textit{Dask}. Este \emph{scheduler}, altamente configurable, permite ejecutar tareas en clústeres distribuidos o en entornos heterogéneos, y podría ofrecer mejoras significativas en escenarios con grandes volúmenes de datos o requerimientos de alta disponibilidad. Su exploración permitiría comprender mejor los límites de escalabilidad del sistema y ajustar dinámicamente la estrategia de ejecución según la demanda computacional.

Por otro lado, el análisis estático del código reveló un \gls{mi} bajo en los módulos encargados de la lógica base de los extractores y del registro de todos los extractores del sistema. Este resultado era esperable dado que la clase abstracta \texttt{Extractor}, responsable del comportamiento general de todos los extractores, aumentó en complejidad con el objetivo de simplificar la implementación de nuevos extractores. A pesar de esto, esta observación señala áreas concretas del sistema donde sigue siendo necesario mejorar la estructura interna del código y facilitar su mantenimiento a largo plazo.

Otra línea de trabajo futuro consiste en extender la funcionalidad de persistencia actualmente limitada a los formatos JSON y YAML. Incorporar compatibilidad con otros formatos de serialización o con bases de datos livianas como \textit{SQLite} podría favorecer la interoperabilidad con otros sistemas y optimizar los flujos de trabajo en distintos contextos de análisis. Siguiendo esta línea de pensamiento, consideramos viable la posibilidad de abstraer el \emph{back end} destinado a la ejecución paralela, dando lugar a la posibilidad de reemplazar \textit{Dask} por otras bibliotecas como \textit{Ray} o \textit{Joblib}. Así, sería posible adaptar la estrategia de paralelización a las necesidades específicas del entorno de ejecución sin modificar la lógica de alto nivel de la aplicación.

Por último, otra funcionalidad que consideramos relevante para el futuro de \gls{feets} como herramienta de análisis es la incorporación de capacidades de visualización para las características extraídas. Proponemos integrar bibliotecas como \textit{Matplotlib} y \textit{Seaborn} para extender la interfaz de los extractores con un nuevo método reimplementable destinado a la generación de gráficos individuales por característica. Esta funcionalidad permitiría también enriquecer el módulo \texttt{Features} para facilitar el análisis visual de características extraídas sobre múltiples \glspl{lc}.


